\section{Reflected cyclic covers}\label{sec:reflection}

The proof begins with a chamber whose double cyclic cover admits a
marked filling and whose odd cyclic covers of degree at least three
do not.  We state the precise chamber data below;
Sections~\ref{sec:chambers}--\ref{sec:formal}
construct them.  This section then passes from marked fillings to
closed manifold models, using the reflection construction and its
filling implication from~\cite{PD}.

\subsection{The reflection input}

We recall the finite version of the reflection construction.  It uses
the framework of Davis~\cite{Davis1983,Davis2008}, with the chamber
and group calculations established in~\cite[Section~3]{PD}.
Let $B$ be a space with a collared boundary $S$, and choose a finite
flag triangulation $\mathcal L$ of $S$.  Here \emph{flag} means that
every set of pairwise adjacent vertices spans a simplex.  For a vertex
$s$ of $\mathcal L$, let $S_s$ be its closed dual block in the
barycentric subdivision.  These blocks form the panels of $S$.
If $S$ is a three-manifold, a nonempty intersection
$\bigcap_{s\in I}S_s$ is a ball of dimension $4-|I|$; such an
intersection occurs exactly when $I$ spans a simplex of $\mathcal L$.

Write $\mathcal V$ for the vertex set, put
$A=(\mathbb Z/2)^{\mathcal V}$, and denote its coordinate vectors by
$e_s$.  For $x\in S$, set $I(x)=\{s:x\in S_s\}$ and
$A_{I(x)}=\langle e_s:s\in I(x)\rangle$.  For $x$ in the interior
of $B$, set $A_{I(x)}=0$.  Define
\begin{equation}\label{refl:finite-reflection}
 \mathcal R_A(B)=(A\times B)/\sim,
 \qquad
 (a,x)\sim(a',x) \Longleftrightarrow\ a-a'\in A_{I(x)}.
\end{equation}
Thus chambers are labeled by $A$, and crossing the panel $S_s$
changes the label by $e_s$.  This is the quotient of the right-angled
Coxeter development by the kernel of the homomorphism sending each
panel generator to its own coordinate in $A$.

The following statement is extracted from the proof in
\cite[Sections~3--5]{PD}.  It includes the precise manifold-to-filling
implication that we require, obtained there by stable realization
and controlled removal of middle homology.

\begin{theorem}[Reflection and marked fillings, \cite{PD}]
\label{refl:input}
Let $D$ be a compact connected oriented smooth four-manifold with
connected boundary $S$ and free fundamental group $H$.  Suppose that
$D$ has the homotopy type of a finite graph with $2k$ two-spheres,
and that its whole middle module $H_2(\widetilde D;\mathbb Z)$ over
$\mathbb ZH$ has a specified
framed immersed hyperbolic basis of $k$ planes, with vanishing
quadratic self-intersections.  Attach three-cells along the basis
spheres to obtain a map of collared pairs
\[
 (D,S)\longrightarrow(P,S)
\]
that is the identity on $S$ and induces the given identification of
fundamental groups.  Assume that $(P,S)$ is an oriented Poincar\'e
pair of dimension four, that $P\simeq BH$, and that
$\pi_1(S)\to H$ is surjective.

For a finite flag triangulation of $S$, let
$Q=\mathcal R_A(P)$ be \eqref{refl:finite-reflection}.  Then $Q$ is a
finite aspherical oriented Poincar\'e complex of dimension four.
If $Q$ has a closed topological four-manifold model, then $S$ admits
a marked graph-type filling: there is a compact oriented topological
four-manifold $V$ with an orientation-preserving boundary identification
$\partial V=S$ and a homotopy
equivalence $V\to P$ whose boundary restriction is homotopic to the
given inclusion $S\to P$.
\end{theorem}

For the first assertion, the relevant reflected-pair statement is
\cite[Proposition~3.2]{PD}, together with its universal-development
argument.  Section~3 of that paper also constructs a proper cocompact
CAT(0) action for the reflected group and obtains the Farrell--Jones
assembly results used in its stable realization argument.  Thus these
are consequences of the stated chamber input, rather than additional
hypotheses on a model.  Under the manifold-model hypothesis,
\cite[Proposition~4.6]{PD} supplies periodically labeled standard
hyperbolic pairs after stabilization, with the prescribed chamber
markings.  The removal argument ends in
\cite[Proposition~5.5]{PD}, which produces the marked filling.  These
arguments use the chamber hypotheses stated in
Theorem~\ref{refl:input}; the particular obstruction to filling in
that paper is applied only after the filling has been produced.
The orientation of a hypothetical model is obtained from its
homotopy equivalence to the oriented Poincar\'e complex $Q$.

\subsection{The parity chamber}

The following proposition isolates the geometric data used in the
rest of the proof.  In particular, it records the lifted chamber
hypotheses needed to apply Theorem~\ref{refl:input} in every cyclic
cover.

\begin{proposition}[Parity chamber data]\label{refl:chamber-data}
There are chamber data $(D,P,S,H)$ satisfying the chamber hypotheses of
Theorem~\ref{refl:input}, with $H$ a finitely generated nonabelian free
group, and a distinguished epimorphism $w\colon H\to\mathbb Z$ with the following
properties.  For $d\geq1$, put $H_d=w^{-1}(d\mathbb Z)$, let
$D_d,P_d$ be the associated connected cyclic covers, and let $S_d$ be
the preimage of $S$.
\begin{enumerate}
\item For every $d$, the boundary $S_d$ is connected, and
$(D_d,P_d,S_d,H_d)$ satisfies all the chamber hypotheses of
Theorem~\ref{refl:input}, with the full framed immersed hyperbolic
basis and its vanishing quadratic self-intersections lifted from $D$.
\item There is a compact oriented topological four-manifold $V$ with
an orientation-preserving boundary identification $\partial V=S_2$
and a homotopy equivalence
\[
 j:(V,S_2)\longrightarrow(P_2,S_2)
\]
equal to the identity on the boundary.
\item For every odd $d\geq3$, there is no compact oriented topological
four-manifold $V_d$ with an orientation-preserving boundary
identification $\partial V_d=S_d$ and a homotopy equivalence
$V_d\to P_d$ whose boundary restriction is homotopic to the inclusion
$S_d\to P_d$.
\end{enumerate}
\end{proposition}

\begin{proof}
Propositions~\ref{ch:homotopy-type} and~\ref{ch:poincare-chamber},
together with Corollary~\ref{ch:boundary-surjectivity}, construct the
base chamber and verify the chamber hypotheses of Theorem~\ref{refl:input}.
The character is the one in \eqref{ch:cover-character}.
Lemma~\ref{ch:cover-data} proves the assertions about every lifted
chamber.  Proposition~\ref{ch:even-filling} gives the relative filling
$j$, and Proposition~\ref{for:odd-nonfilling} gives the last assertion,
using the marking in Definition~\ref{ch:marked-filling}.  These results
are proved in Sections~\ref{sec:chambers}--\ref{sec:formal} from the
explicitly stated inputs of~\cite{MT,PD}.
\end{proof}

\subsection{Comparing reflected cyclic covers}

Fix the chamber of Proposition~\ref{refl:chamber-data}.  Choose any
finite flag triangulation $\mathcal L$
of $S$, with vertex set $\mathcal V$, and put
\[
 \mathsf A_0=(\mathbb Z/2)^{\mathcal V},\qquad
 Q^0=\mathcal R_{\mathsf A_0}(P),\qquad G^0=\pi_1(Q^0).
\]
There is a folding map
\begin{equation}\label{refl:folding}
 r^0\colon Q^0\longrightarrow P,
 \qquad r^0([a,x])=x.
\end{equation}
Its restriction to every chamber is the identity.  Thus $r^0_*$ is a
split epimorphism onto $H$.  Compose it with the epimorphism
$w\colon H\to\mathbb Z$ in \eqref{ch:cover-character}, and use
$w$ also for this epimorphism $G^0\to\mathbb Z$.
For a positive integer $d$, let $Q^0_d\to Q^0$ denote its connected
degree-$d$ cyclic cover and put
$G^0_d=\pi_1(Q^0_d)=w^{-1}(d\mathbb Z)$.

The pullback of $P_d\to P$ along \eqref{refl:folding} is $Q^0_d$.
In particular, $Q^0_d$ has the following concrete chamber description:
use one copy of $P_d$ for each label in $\mathsf A_0$, and cross every lift of
$S_s$ by adding the original coordinate $e_s$.  This pullback is
connected because the composite $G^0\to\mathbb Z/d$ is surjective.
The boundary of each chamber carries the lifted triangulation
$\mathcal L_d$ of $S_d$.

Theorem~\ref{refl:input} uses an independent coordinate for every panel
of its chamber boundary, whereas $Q^0_d$ reuses the downstairs
coordinates.  The next lemma separates the lifted panel labels by
passing to a finite cover, to which the theorem applies directly.

\begin{lemma}\label{refl:cover-comparison}
The triangulation $\mathcal L_d$ is flag.  Let
$\mathsf A_d=(\mathbb Z/2)^{\mathcal V_d}$, where $\mathcal V_d$ is its
vertex set, and distinguish all its lifted panels in the reflection
\[
 \widehat Q_d=\mathcal R_{\mathsf A_d}(P_d).
\]
There is a connected finite covering $\widehat Q_d\to Q^0_d$.
\end{lemma}

\begin{proof}
A clique of vertices in $\mathcal L_d$ projects to distinct pairwise
adjacent vertices of $\mathcal L$.  Flagness downstairs supplies a
simplex containing their images.  Lift that simplex from one vertex
of the clique.  The clique edges incident to that vertex project to
edges of this simplex; uniqueness of their lifts places all vertices
of the clique in the chosen lifted simplex.  Thus
$\mathcal L_d$ is flag.

Define a surjective homomorphism
\[
 \pi_d\colon \mathsf A_d\longrightarrow\mathsf A_0
\]
by sending the coordinate of a lifted panel to the coordinate of its
downstairs panel.  In the chamber description of $Q^0_d$ above, the
formula $[a,x]\mapsto[\pi_d(a),x]$ is well-defined and surjective.
At a point belonging to a set of meeting lifted panels, those panels
project to distinct panels downstairs.  The homomorphism $\pi_d$
therefore restricts to an isomorphism on their coordinate subgroups.
The sector identifications in \eqref{refl:finite-reflection} consequently
give covering charts at this point; in chamber interiors the assertion
is immediate.  The map is finite because $\mathsf A_d$ is finite.  Finally,
$P_d$ is connected, and each coordinate generator of $\mathsf A_d$ is realized
by crossing its nonempty panel.  All its chambers are therefore in
one connected component.
\end{proof}

\begin{proposition}\label{refl:odd-models}
For every odd integer $d\geq3$, the space $Q^0_d$ has no closed
topological four-manifold model.
\end{proposition}

\begin{proof}
Suppose that $Q^0_d$ has such a model.  The finite cover of
Lemma~\ref{refl:cover-comparison} determines, through the model
equivalence, a finite cover of that closed manifold.  Lifting the
equivalence gives a closed topological four-manifold model for
$\widehat Q_d$.

Proposition~\ref{refl:chamber-data} verifies all the hypotheses of
Theorem~\ref{refl:input} for $(D_d,P_d,S_d,H_d)$, with the lifted
hyperbolic basis and boundary marking.  The theorem therefore gives a
marked filling of $S_d$, contradicting the last assertion of
Proposition~\ref{refl:chamber-data}.
\end{proof}

\subsection{The triangulation used in the construction}

For the rest of the paper, choose $\mathcal L$ to be a finite
flag-no-square PL triangulation of $S$ with at least $5000$ vertices.
Here \emph{flag-no-square} means flag with no induced four-cycle in
the one-skeleton.  Such a triangulation exists by
\cite[Proposition~2.13 and Corollary~2.14]{PrzytyckiSwiatkowski2009}:
start with a finite PL triangulation, subdivide until it has at least
$5000$ vertices, and then use their flag-no-square subdivision, which
retains the old vertices.  We keep the notation $\mathcal V$,
$\mathsf A_0$, $Q^0$, and $Q^0_d$ for this choice.

For any finite PL triangulation used below, we make compatible cellular
choices for the later chamber gluings.  Give $S$ the subdivision in which every dual panel is a
subcomplex, and extend this finite triangulation over a collar and
then over the smooth manifold $D$.  Cellular approximation of the
finitely many interior three-cell attaching maps gives a finite CW
pair $(P,S)$ in the same homotopy type relative to $S$.  These choices
pass to finite covers.  In particular, every union of panels is a
boundary subcomplex and its inclusion in a collared chamber is a
cofibration.  Replacing the chamber by this relative CW model changes
no reflected homotopy type: adjoining the finitely many chambers in
turn and using the gluing lemma for homotopy pushouts proves this from
the relative homotopy equivalence.  We use these models without
changing notation.
